\section{Service Architecture Used}
\label{sec:arch}

\subsection{Services and data flows}

The application is a temperature-monitoring pipeline written in Python (Flask). Each service is a
separate process, image and Kubernetes Deployment (Table~\ref{tab:services}). A sensor posts a
reading every 5~s to an ingestion instance, which validates it and forwards it to storage; the
analytics service reads through the storage API, never the database file. Every service finds its
dependencies through the registry at each call, so a replaced instance is used as soon as it
registers.

\begin{table}[ht]
\centering\small
\caption{Services, ports and their own Prometheus metrics (besides the common ones).}
\label{tab:services}
\begin{tabularx}{\textwidth}{@{}l l Y Y@{}}
\toprule
Service & Port & Role & Own metrics \\
\midrule
registry  & 5050 & registration with heartbeat, discovery by location, round-robin, health checks, \code{/tree} & \metric{registry\_instances} (per service and location), registrations, removals, failed health checks \\
ingestion & 5001 & validate readings, forward to storage (502 if storage fails) & \metric{ingestion\_readings\_total\{result\}} (accepted, invalid, storage\_error) \\
storage   & 5002 & SQLite persistence and read API (single writer) & stored, rejected, \metric{storage\_rows} \\
analytics & 5003 & average/min/max over stored readings & \metric{analytics\_storage\_errors\_total} \\
sensor    & 8000 & posts a reading every 5~s (metrics only) & \metric{sensor\_readings\_sent\_total\{result\}} \\
\bottomrule
\end{tabularx}
\end{table}

All HTTP services share one instrumentation module that exports
\metric{http\_requests\_total}, a latency histogram, the number of in-flight requests, and
\metric{service\_info}, a constant series that carries the instance's ID and location. Labels never
contain client-controlled values (for example the sensor's device ID), so a client cannot inflate
the number of time series. The same module writes one JSON access-log line per request
(\code{ts}, \code{service}, \code{id}, \code{srcip}, method, path, status, latency), which is what the
Wazuh agents read.

\subsection{Deployment: physical and logical location}

Figure~\ref{fig:layers} shows where every part runs, from the one real machine down to the
programs; Figure~\ref{fig:arch} adds the data flows and trust boundaries.

\begin{figure}[ht]
\centering
\resizebox{\textwidth}{!}{% Figure: where every part runs, from the one real machine down to the pods.
% Source facts: FACTS.md sections 1, 3, 4 (k3s/setup-cluster.sh, k8s/ssle.yaml, README.md).
\begin{tikzpicture}[
  font=\small,
  title/.style={font=\small\bfseries, anchor=north west, inner sep=0pt},
  sub/.style={font=\footnotesize, anchor=north west, inner sep=0pt, text=black!65},
  item/.style={font=\footnotesize, anchor=north west, inner sep=0pt},
]
% one real machine
\draw[rounded corners=10pt, fill=gray!8, draw=gray!60] (0,0) rectangle (15.6,10.2);
\node[title] at (0.4,9.85) {One Kali VM: the only real machine};
\node[sub]   at (0.4,9.35) {every box inside is an Incus container};

% k3s cluster
\draw[rounded corners=8pt, fill=violet!7, draw=violet!55] (0.4,2.3) rectangle (9.9,8.8);
\node[title, text=violet!70!black] at (0.75,8.5) {k3s cluster (Kubernetes)};
\node[sub] at (0.75,8.0) {2 nodes = 2 simulated data centres};

% node dc1
\draw[rounded corners=6pt, fill=teal!8, draw=teal!60] (0.75,2.65) rectangle (5.05,7.55);
\node[title, text=teal!60!black] at (1.05,7.3) {k3s-master};
\node[sub] at (1.05,6.8) {node label: dc1, rack1};
\foreach \t [count=\i] in {registry, storage (SQLite), ingestion-dc1, analytics-dc1, sensor-dc1, Prometheus}
  \node[item] at (1.05,6.75-0.55*\i) {pod: \t};

% node dc2
\draw[rounded corners=6pt, fill=teal!8, draw=teal!60] (5.3,2.65) rectangle (9.55,7.55);
\node[title, text=teal!60!black] at (5.6,7.3) {k3s-node1};
\node[sub] at (5.6,6.8) {node label: dc2, rack1};
\foreach \t [count=\i] in {ingestion-dc2, analytics-dc2, sensor-dc2}
  \node[item] at (5.6,6.75-0.55*\i) {pod: \t};

% outside the cluster
\draw[rounded corners=8pt, fill=orange!8, draw=orange!60] (10.25,2.3) rectangle (15.2,8.8);
\node[title, text=orange!60!black] at (10.6,8.5) {Outside the cluster};
\node[sub] at (10.6,8.0) {shared by both sites};
\foreach \t [count=\i] in {Wazuh manager (SIEM), central Prometheus, mail server, DNS server, docker: builds images}
  \node[item] at (10.6,7.6-0.65*\i) {\t};
\node[item] at (10.6,3.95) {\phantom{docker: }+ private registry};

% images are pulled by the nodes
\draw[-{Stealth[length=2mm]}, thick, orange!70!black] (12.0,3.45) |- (9.55,3.0)
  node[pos=0.75, above, font=\footnotesize] {images};

% reading key
\node[item] at (0.4,1.85) {\textbf{Physical level} = the node (which machine it runs on)};
\node[item] at (0.4,1.3)  {\textbf{Logical level} = the pod (which program instance it is)};
\node[item] at (0.4,0.7)  {ID: \texttt{ingestion-service-ingestion-dc2-<hash>.k3s-node1.rack1.dc2.centro.pt}};
\end{tikzpicture}
}
\caption{Where every part runs. The only physical machine is one VM; every box inside it is an
Incus container. Two of them form the k3s cluster, and each node is labelled as a data centre. The
node is the physical level of an ID and the pod the logical level.}
\label{fig:layers}
\end{figure}

Two k3s nodes run as privileged Incus containers on
one VM; each node \emph{is} a data centre: \code{k3s-master} is \code{dc1} and \code{k3s-node1} is
\code{dc2}, both in \code{rack1}, region \code{centro}, country \code{pt} (node labels
\code{ssle.dc}, \code{ssle.rack}). Ingestion, analytics and the sensor run once per location, as one
Deployment each pinned with a \code{nodeSelector}; the registry, storage (one replica, persistent
volume) and the in-cluster Prometheus run in dc1. Images are built in a separate container and
pulled from a private registry, tagged with the git commit. The Wazuh manager, the central
Prometheus, the mail server and the DNS server run outside the cluster on the same host.

\begin{figure}[ht]
\centering
\resizebox{\textwidth}{!}{% Figure 1: deployment on k3s (main platform) with trust boundaries.
% Source facts: FACTS.md sections 2-6 (k8s/ssle.yaml, k3s/README.md, k3s/setup-cluster.sh).
\begin{tikzpicture}[
  font=\scriptsize,
  svc/.style={draw, rounded corners=2pt, fill=blue!8, minimum width=24mm, minimum height=8mm, align=center},
  int/.style={svc, fill=gray!15},
  sec/.style={svc, fill=orange!18},
  ext/.style={svc, fill=red!10},
  flow/.style={-{Stealth[length=1.8mm]}, semithick},
  mon/.style={-{Stealth[length=1.8mm]}, dashed, thick, gray!60!black},
  bound/.style={draw, dashed, rounded corners=4pt, inner sep=6pt},
  lbl/.style={font=\scriptsize\bfseries, anchor=south west, inner sep=1pt}
]
% ---------------- node k3s-master = dc1
\node[svc] (sen1) at (1.4,3.0)  {sensor-dc1};
\node[svc] (ing1) at (1.4,1.6)  {ingestion-dc1\\NodePort 30001};
\node[svc] (ana1) at (4.6,1.6)  {analytics-dc1\\NodePort 30003};
\node[int] (reg)  at (1.4,0.2)  {registry\\\textit{ClusterIP 5050}};
\node[int] (sto)  at (4.6,0.2)  {storage (SQLite)\\\textit{ClusterIP 5002}};
\node[sec] (ag1)  at (1.4,-1.2) {Wazuh agent\\+ k3s-block};
\node[int] (pin)  at (4.6,-1.2) {Prometheus\\(in-cluster) 30090};
\node[bound, fit=(sen1)(ana1)(ag1)(pin)] (n1) {};
\node[lbl] at (n1.north west) {k3s-master = dc1 / rack1};

% ---------------- node k3s-node1 = dc2
\node[svc] (sen2) at (8.6,3.0)  {sensor-dc2};
\node[svc] (ing2) at (8.6,1.6)  {ingestion-dc2\\NodePort 30001};
\node[svc] (ana2) at (11.4,1.6) {analytics-dc2\\NodePort 30003};
\node[sec] (ag2)  at (10.0,-1.2) {Wazuh agent\\+ k3s-block};
\node[bound, fit=(sen2)(ana2)(ag2)] (n2) {};
\node[lbl] at (n2.north west) {k3s-node1 = dc2 / rack1};

\begin{scope}[on background layer]
  \node[draw, thick, rounded corners=6pt, fill=blue!2, inner sep=13pt, fit=(n1)(n2)] (cl) {};
\end{scope}
\node[font=\scriptsize\bfseries, anchor=south] at ([yshift=6pt]cl.north)
  {k3s cluster: pod network, region centro, country pt};

% ---------------- outside the cluster (same Incus host, lab network)
\node[sec] (wz)   at (0.2,-4.4)  {Wazuh manager\\rules 1002xx, AR};
\node[int] (pc)   at (3.5,-4.4)  {central Prometheus\\(federation)};
\node[ext] (mail) at (6.85,-4.4)  {mail server\\Postfix + SpamAssassin};
\node[ext] (dns)  at (10.2,-4.4)  {DNS server (dnsmasq)\\\code{*.ssle.local}};
\node[int] (dreg) at (13.5,-4.4) {image registry\\:5000};
\node[ext] (cli)  at (15.2,1.6)  {clients /\\attackers};

% ---------------- data flows
\draw[flow] (sen1) -- (ing1);
\draw[flow] (ing1) -- (sto) node[midway, above, sloped] {store};
\draw[flow] (ana1) -- (sto) node[midway, right] {read};
\draw[flow, gray] (ing1) -- (reg) node[midway, left] {lookup};
\draw[flow] (sen2) -- (ing2);
\draw[flow] (ing2.south) |- (sto.east) node[pos=0.75, below] {store (cross-dc)};
\draw[flow] (cli) -- (ana2);
\node[font=\tiny, align=center, below=1mm of cli] {via each location's\\NodePorts\\(\code{externalTrafficPolicy:}\\\code{Local})};

% ---------------- monitoring and security
\draw[mon] (pc.north) -- (pc.north |- pin.south) node[pos=0.35, right] {federate \code{location:*}};
\draw[mon] (ag1.south) -- ([xshift=10mm]wz.north);
\draw[mon] (ag2.south) -- ++(0,-1.05) -| ([xshift=-6mm]wz.north);
\draw[mon] (dns.south) -- ++(0,-0.45) -| (wz.south);
\draw[dashed, thick, gray!60!black] (mail.south) -- ++(0,-0.45);
\draw[mon] (dreg.north) -- (dreg.north |- cl.south) node[pos=0.5, right] {images};
\end{tikzpicture}
}
\caption{Main deployment on k3s: each node is a data centre. Grey services are internal
(ClusterIP), blue ones are published on NodePorts, dashed lines are monitoring and security
traffic.}
\label{fig:arch}
\end{figure}

The baseline deployment uses one Incus container per service with a systemd unit; the location
is given by environment variables. All tests and evidence exist for both platforms
(Section~\ref{sec:eval}), and the comparison is itself informative: three of the problems reported
in Section~\ref{sec:threats} only appeared on Kubernetes.

\subsection{Identification and registration of elements}
\label{sec:naming}

Every instance builds its ID from five location levels, most specific first, like a DNS
name~\cite{rfc1034}:
\[
\underbrace{\code{ingestion-service-ingestion-dc2-<hash>}}_{\text{service-instance (logical)}}.
\underbrace{\code{k3s-node1}}_{\text{node (physical)}}.\code{rack1}.\code{dc2}.\code{centro}.\code{pt}
\]
The node level is the machine that runs the element and the instance level is the logical element
on it: on Kubernetes they come from \code{spec.nodeName} and the pod name (downward API); on Incus
from the host name and an instance number (Figure~\ref{fig:idtree}). Dots inside a level are
replaced, so a label can never fake another level.

\begin{figure}[ht]
\centering
\resizebox{0.85\textwidth}{!}{% Figure 2: naming hierarchy and the k3s mapping. Source: common.py, k8s/ssle.yaml, FACTS.md section 3.
\begin{forest}
  for tree={font=\scriptsize, draw, rounded corners=2pt, align=center, l sep=4mm, s sep=2mm,
            edge={-{Stealth[length=1.3mm]}}, fill=blue!5}
  [{country\\\code{pt}}
    [{region\\\code{centro}}
      [{dc\\\code{dc1}\\(k3s node label)}
        [{rack\\\code{rack1}}
          [{node\\\code{k3s-master}\\(\code{spec.nodeName})}
            [{instance\\\code{storage-<pod>}}, fill=orange!15]
            [{instance\\\code{ingestion-dc1-<pod>}}, fill=orange!15]
          ]
        ]
      ]
      [{dc\\\code{dc2}}
        [{rack\\\code{rack1}}
          [{node\\\code{k3s-node1}}
            [{instance\\\code{ingestion-dc2-<pod>}}, fill=orange!15]
          ]
        ]
      ]
    ]
  ]
\end{forest}
}
\caption{Location hierarchy of the IDs. Orange leaves are the logical elements (pods); the node
level is the physical machine.}
\label{fig:idtree}
\end{figure}

Registration is decentralised: each instance posts \code{\{name, url, id, location\}} to the
registry every 10~s. The call is an idempotent upsert, so it doubles as a heartbeat and recovers
automatically when the registry restarts. On Kubernetes each pod registers its own pod IP; with the
Service address all replicas collapsed into one entry, a problem found and fixed during development. The registry checks
every instance's \code{/health} every 10~s and removes it after two consecutive failures. Queries
can be filtered by any level (\code{/services?dc=dc2}), and \code{/tree} returns instance counts per
country, region, dc, rack and node.

\subsection{Geographic location and delegation of responsibilities}

Location is used, not only recorded:
\begin{itemize}
  \item \textbf{Discovery prefers the closest instance.} A service first asks for instances in its own
        location up to a configurable scope (default: same dc), then any location, then a static
        fallback. The dc2 sensor therefore feeds the dc2 ingestion pod, as verified
        (Section~\ref{sec:eval}).
  \item \textbf{Each location serves its own clients.} The NodePorts use
        \code{externalTrafficPolicy: Local}, so a node only forwards to its own pods and keeps the
        client's source address~\cite{k8s_sourceip}, which the SIEM needs to block it.
  \item \textbf{Monitoring is delegated per site.} The in-cluster Prometheus scrapes the local pods;
        the central one only federates per-location aggregates and alerts.
  \item \textbf{Security enforcement is delegated per node.} One Wazuh agent per k3s node reads the
        logs of every pod on that node and executes the active responses locally; the manager only
        decides.
\end{itemize}

\subsection{Monitoring by location with optimised bandwidth}

Recording rules join every metric with \metric{service\_info} and sum it by country, region and dc:
\metric{location:http\_requests:rate1m}, \metric{location:http\_errors:rate1m},
\metric{location:http\_request\_duration\_seconds:p95} and \metric{location:instances:count}. The
central Prometheus federates only \code{\{\_\_name\_\_=\textasciitilde"location:.+"\}} (plus
\code{ALERTS}), following the hierarchical federation pattern~\cite{prometheus_federation}. Its
traffic grows with the number of locations, not with the number of pods; the measured saving is in
Section~\ref{sec:eval}.

\subsection{Locating the source of an event}

Because the same ID appears in the registry, in \metric{service\_info}, in every access-log line
and, at node level, in the Wazuh agent name (\code{k3s-node-01.k3s-node1.rack1.dc2.centro.pt}), any
event can be placed without a lookup table: a Wazuh alert names the node and carries the pod's full
ID in \code{data.id}; a Prometheus alert can be joined with \metric{service\_info} to get the
location. For example, the alerts of the DoS test identify both the attacker (\code{srcip}) and the
attacked pod in dc1 (Section~\ref{sec:threats}).

\subsection{Exposure and trust boundaries}

Only ingestion and analytics are reachable from outside the cluster (NodePorts 30001/30003). The
registry and storage are ClusterIP services: the registry has no authentication, so it must never
be exposed. Pods have CPU and memory limits (500m, 256~MiB) and readiness and liveness probes, which
both restart failed pods and cap what a flood can take from a node. The trust boundaries are
therefore: the lab network (clients, attackers, mail, DNS), the node (NodePorts, agent, firewall),
and the pod network (internal services, trusted by the SIEM rules, see
Section~\ref{sec:threats}). The k3s nodes run as \emph{privileged} containers, which is part of the
attack surface analysed next.
